\documentclass{article}
\usepackage[utf8]{inputenc}
\usepackage[turkish]{babel}
\usepackage{booktabs}
\usepackage{amsmath}
\usepackage[hyphens]{url}
\usepackage{hyperref}
\usepackage{geometry}
\geometry{a4paper, margin=2.5cm}
\sloppy

\title{\textbf{Hepiniz Kanye West Konserine Gittiniz}}
\date{}

\begin{document}
\maketitle

Bir süredir sosyal medyada Kanye West konseri hakkında süregelen tartışmayı takip ediyorum. Bu konuda benim yaklaşımım sanatçıların kişilikleri ve sanatlarını ayırmak, ancak finansal destek sağlamamak yönündeydi. Sanatçılar streaming platformlardan pek gelir elde etmediğinden, stream etmenin kayda değer bir etkisi olmadığını; merch almak ya da konsere gitmek gibi eylemlerin ise problematik olabileceğini düşünüyordum.

Asıl rahatsız olduğum şey ise şu: Kanye West konserine gidenlere ya da onu dinleyenlere tepki gösterenlerin, başka problematik kişilere yaptıkları finansal katkıları görmezden gelmesi.

Kabul edelim ki hiçbirimiz bu tür katkılardan tamamen kaçınamıyoruz. Ama bir yandan finansal destek verip bunu görmezden gelirken, öte yandan başkalarını bu yüzden linçlemek kendi içinde tutarsız ve problematik. Çünkü dediğim gibi, hepimiz insanız.

Uyarmanın ya da sorunlardan bahsetmenin yanlış olduğunu düşünmüyorum. Ancak sosyal medyada ``arkadaşlığımı keserim''den bazı ünlüleri linçlemeye uzanan bu tablo beni rahatsız etti. Bu yüzden, Kanye West konserine tepki gösterenlerin yaptığı başka bir şeyin farklı bir problematik kişiye ne kadar finansal destek sağladığına bakmak ve iddialarımı sayısal olarak test etmek istedim: Instagram'da geçirdiğiniz zamanla Mark Zuckerberg'e ne kadar para kazandırdığınızı hesapladım.

\section*{Zuckerberg Neden Problematik?}

Rakamları sunmadan önce, neden Zuckerberg'den bahsettiğimi kısaca açıklamak istiyorum.

\begin{itemize}

\item \textbf{2016 Seçimleri:} Cambridge Analytica, Facebook'un API'sindeki bir açığı kullanarak milyonlarca kullanıcının verisini ele geçirdi ve bu veriyi Trump'ın 2016 kampanyasında seçmenleri hedeflemek için kullandı. Her iki kampanya da bu verileri kullandıklarını resmi olarak kabul etmedi, dolayısıyla gerçek etkisi tartışmalı.\textsuperscript{[1]}

\item \textbf{Fact Checking Sona Erdi:} Zuckerberg, Ocak 2025'te Facebook ve Instagram'daki bağımsız fact checking programını sonlandırarak yerine ``community notes'' sistemini getirdi ve son seçimi ``ifade özgürlüğü açısından kültürel bir kırılma noktası'' olarak tanımladı.\textsuperscript{[2]}

\item \textbf{DEI Programları Kaldırıldı:} Meta, Ocak 2025'te çeşitlilik, eşitlik ve kapsayıcılık (DEI) ekibini feshederek ilgili tüm programları sonlandırdı; işe alım ve tedarikçi çeşitliliği pratiklerini de değiştirdi. Bu adım, Trump'ın göreve dönüşüyle ve daha geniş bir içerik moderasyonu geri adımıyla eş zamanlı gerçekleşti.\textsuperscript{[3]}

\item \textbf{LGBTİ+ İçerik Baskısı:} Meta'nın Instagram'da \#lesbian, \#bisexual, \#gay, \#trans gibi etiketleri varsayılan olarak gençler için açık olan ``hassas içerik'' filtresi kapsamına aldığı, dolayısıyla LGBTİ+ içeriğe erişimi fiilen engellediği ileri sürüldü.\textsuperscript{[4]} Ayrıca Meta'nın güncellenmiş nefret söylemi politikası, kullanıcıların LGBTİ+ bireyleri ``akıl hastası'' olarak nitelendirmesine açıkça olanak tanıyor; Zuckerberg bu kısıtlamaları ``ana akım söylemle bağdaşmayan'' uygulamalar olarak çerçeveledi.\textsuperscript{[5]} 2025'in sonuna gelindiğinde, 50'den fazla LGBTİ+ ve üreme sağlığı kuruluşunun Meta platformlarındaki hesapları askıya alınmış ya da gölge yasakla karşılaşmıştı; belgelenen vakalar 2024'e kıyasla ikiye katlandı.\textsuperscript{[6]}

\item \textbf{Lobi Faaliyetleri:} Meta, 2024'te federal lobi için rekor düzeyde 24,4 milyon dolar harcadı; bu rakam 2023'e göre yüzde 27 artış anlamına geliyor ve 2009'dan bu yana görülen en yüksek lobi harcaması. Şirket bu süreçte 65 lobici istihdam etti; yani ABD Kongresi'nin her sekiz üyesine bir lobici düşüyor.\textsuperscript{[7]}

\item \textbf{Peter Thiel:} Peter Thiel, 2005'te erken dönem yatırımcı sıfatıyla Facebook'un yönetim kuruluna katıldı ve yaklaşık yirmi yıl boyunca Zuckerberg'in yakın danışmanı olarak görev yaptı; Zuckerberg'in iş anlayışını şekillendirmede belirleyici bir rol oynadığı ifade edildi. Thiel, 2022'de Trump'ın siyasi gündemini desteklemek amacıyla yönetim kurulundan ayrıldı ve Senato seçimlerinde J.D. Vance ile Blake Masters'ı destekledi.\textsuperscript{[8]}

\end{itemize}

``E başka ne yapalım?'' diyecekseniz, Bluesky, Substack, Mastodon, hatta Reddit bile Meta ve X'ten kat kat daha az problematik seçenekler.

\section*{Hesaplamalar}

\subsection*{Bir Türk Kullanıcısının Meta İçin Yıllık Değeri}

Meta'nın kullanıcı başına geliri (ARPU) şu formülle hesaplanıyor:

\[
\text{ARPU} = \frac{\text{Toplam Gelir}}{\text{Aylık Aktif Kullanıcı Sayısı}}
\]

Bu yöntemin bilerek tercih edilmesinin nedeni şu: Meta'nın geliri doğrudan her kullanıcıdan ne kadar zaman çaldığına değil, kaç kullanıcısı olduğuna bağlı. Daha fazla kullanıcı, daha büyük bir reklam havuzu ve daha fazla veri demek; bu da daha yüksek toplam gelir anlamına geliyor. Dolayısıyla ``platformdasın'' demek, ``bu rakama katkıda bulunuyorsun'' demektir. Basitlik adına herkesi ortalama kullanıcı olarak aldım.

Meta, Türk kullanıcılarını kendi raporlamasında \textbf{Avrupa} segmentine dahil ediyor\textsuperscript{[9]}; ancak Türkiye'nin reklam pazarı Batı Avrupa'ya kıyasla daha az gelişmiş. Bu nedenle üç senaryo sunuyorum:

\begin{table}[h]
\centering
\begin{tabular}{lcc}
\toprule
\textbf{Senaryo} & \textbf{Yıllık ARPU (USD)} & \textbf{Türkiye için uygunluk} \\
\midrule
Avrupa (üst sınır)           & \$68,12 & Resmi segment \\
Küresel ortalama             & \$57,03 & Orta tahmin \\
Asya-Pasifik (alt sınır)     & \$21,28 & Pazar olgunluğu benzer \\
\midrule
Türkiye tahmini (orta nokta) & \$39,16 & $(57{,}03 + 21{,}28) / 2$ \\
\bottomrule
\end{tabular}
\caption{Meta ARPU senaryoları\textsuperscript{[10],[11]}}
\end{table}

Yani Instagram ya da Facebook'ta hesabınız varsa, yılda yaklaşık \textbf{\$21--\$68} arasında Zuckerberg'e kazandırıyorsunuz. En muhafazakâr tahminle \textbf{\$39}.

\subsection*{Ayda 50 Dinleme}

Hesaplama için aylık 50 dinleme esas alınmıştır. Bu rakam, Spotify'ın kendi dinleyici segmentasyonuna göre ``süper dinleyici'' eşiğinin (ayda 15+ intentional stream) belirgin biçimde üzerinde, yani bir sanatçının en sadık hayranlarının davranışını temsil eder.\textsuperscript{[12]} Ortalama bir kullanıcı haftada 40 farklı sanatçı dinlediğinden,\textsuperscript{[13]} tek bir sanatçıya ayda 50 stream oldukça bonkör bir tahmindir. Gerçek rakamınız muhtemelen daha düşük --- bu da aşağıdaki sonuçları daha da vahim hâle getirir.

Türkiye'de premium plan abonelik fiyatları: Spotify aylık 99₺ ($\approx$\$2,16), Apple Music 59,99₺ ($\approx$\$1,31), YouTube Music 79,99₺ ($\approx$\$1,74).\textsuperscript{[14],[15]} Her üç platform da pro-rata (havuz) modeli kullanıyor; abonelik fiyatı ne kadar düşükse sanatçıya giden ödeme de o kadar düşük oluyor.

Mevcut Türkiye'ye özgü veriler ve abonelik fiyatlarıyla orantılı ölçekleme kullanılarak yapılan tahminler:

\begin{table}[h]
\centering
\begin{tabular}{lccc}
\toprule
\textbf{Platform} & \textbf{Stream başı (TR)} & \textbf{50 dinleme/ay} & \textbf{Yıllık} \\
\midrule
Spotify       & \$0,001--\$0,002   & \$0,05--\$0,10   & \$0,60--\$1,20 \\
Apple Music   & \$0,0012--\$0,0017 & \$0,06--\$0,085  & \$0,72--\$1,02 \\
YouTube Music & $\approx$\$0,00069 & $\approx$\$0,035 & $\approx$\$0,41 \\
\bottomrule
\end{tabular}
\caption{Türkiye'den ayda 50 stream başına sanatçıya giden ödeme
(Haziran 2026 kuru: 1\$ $\approx$ 46₺)\textsuperscript{[14],[15],[16]}}
\end{table}

Üç platformun birbirine bu kadar yakın çıkması, pro-rata modelin bir sonucu: abonelik fiyatı değişse de havuzun toplam büyüklüğü ve platform politikaları nihai oranı belirliyor. YouTube Music'in geride kalmasının temel nedeni ise reklam destekli ücretsiz katmanın havuzu seyreltmesi.

\subsection*{Kanye West İstanbul Konseri --- 30 Mayıs 2026}

Konser, Atatürk Olimpiyat Stadyumu'nda gerçekleşti. Resmi bilet fiyatları\textsuperscript{[17],[18]}:

\begin{table}[h]
\centering
\begin{tabular}{lcc}
\toprule
\textbf{Kategori} & \textbf{TL} & \textbf{USD ($\approx$46 TL/\$)} \\
\midrule
VIP              & ₺100.000 & $\approx$\$2.174 \\
Pit (sahne önü)  & ₺65.000  & $\approx$\$1.413 \\
Diamond          & ₺35.000  & $\approx$\$761  \\
Batı Alt Premium & ₺23.400  & $\approx$\$509  \\
Golden           & ₺18.000  & $\approx$\$391  \\
Silver           & ₺12.600  & $\approx$\$274  \\
Kuzey / Güney    & ₺7.200   & $\approx$\$157  \\
\bottomrule
\end{tabular}
\caption{Ye Live in İstanbul bilet kategorileri ve fiyatları\textsuperscript{[17],[18]}}
\end{table}

\subsection*{Karşılaştırma}

\begin{table}[h]
\centering
\begin{tabular}{lc}
\toprule
\textbf{Eylem} & \textbf{Yıllık finansal katkı (USD)} \\
\midrule
Instagram/Facebook kullanımı (Meta'ya)               & \$21--\$68 \\
Spotify'da ayda 50 stream (sanatçıya)                & \$0,60--\$1,20 \\
Apple Music'te ayda 50 stream (sanatçıya)            & \$0,72--\$1,02 \\
YouTube Music'te ayda 50 stream (sanatçıya)          & $\approx$\$0,41 \\
Kanye West İstanbul --- en ucuz bilet (tek seferlik) & $\approx$\$157 \\
\bottomrule
\end{tabular}
\caption{Farklı eylemlerin yarattığı finansal katkı}
\end{table}

\section*{Özetle}

Meseleyi somutlaştırmak gerekirse: en ucuz konser bileti (\$157) değerinde finansal katkıyı Zuckerberg'e yapmak için Instagram'ı yaklaşık \textbf{2--7 yıl} kullanmak yeterli. Kanye'yi Spotify'da ayda 50 kez dinleyerek ona aynı miktarı kazandırmak ise \textbf{131--262 yıl} alıyor.

Yani, bunu okuyan herkesin yaklaşık \textbf{2--7 yıl} civarı Instagram kullandığını düşünürsek:

Hayırlı olsun, hepiniz bir Kanye West konserine gittiniz.

% ---------------------------------------------------------------
\section*{Kaynaklar}

\begin{enumerate}

% --- Meta / Zuckerberg ---
\item TechCrunch --- \textit{Facebook Suspends Cambridge Analytica} (2018).
  \url{https://techcrunch.com/2018/03/16/facebook-suspends-cambridge-analytica-the-data-analysis-firm-that-worked-for-the-trump-campaign}
\item Washington Post --- \textit{Meta Ends Fact-Checking} (2025).
  \url{https://www.washingtonpost.com/technology/2025/01/07/meta-factchecking-zuckerberg/}
\item CNN --- \textit{Meta Ends DEI Programs} (2025).
  \url{https://www.cnn.com/2025/01/10/tech/meta-ends-dei-programs/index.html}
\item Metro Weekly --- \textit{Meta Accused of Hiding LGBTQ+ Content} (2025).
  \url{https://www.metroweekly.com/2025/01/meta-accused-of-hiding-lgbtq-content-on-instagram/}
\item Rough Draft Atlanta --- \textit{Meta Content Moderation Policies} (2025).
  \url{https://roughdraftatlanta.com/2025/01/14/meta-content-moderation-policies/}
\item LGBTQ Nation --- \textit{Meta Accused of Banning LGBTQ+ Accounts} (2025).
  \url{https://www.lgbtqnation.com/2025/12/meta-accused-of-banning-lgbtq-accounts-in-one-of-its-biggest-waves-of-censorship-ever/}
\item Issue One --- \textit{Big Tech Spent Record Sums on Lobbying} (2025).
  \url{https://issueone.org/articles/big-tech-spent-record-sums-on-lobbying-last-year/}
\item Fortune --- \textit{Peter Thiel Leaving Meta} (2022).
  \url{https://fortune.com/2022/02/07/peter-thiel-leaving-meta-trump-agenda/}

% --- Meta ARPU ---
\item Stock Dividend Screener --- \textit{Meta Revenue Breakdown by Region} (2026).
  \url{https://stockdividendscreener.com/information-technology/meta-revenue-breakdown-by-region-and-user-geography/}
\item Statista --- \textit{Meta ARPU Worldwide 2025}.
  \url{https://www.statista.com/statistics/234056/facebooks-average-advertising-revenue-per-user/}
\item BusinessStats --- \textit{Facebook ARPU by Region}.
  \url{https://businesstats.com/facebooks-average-advertising-revenue-per-user/}

% --- 50 dinleme gerekçesi ---
\item Spotify for Artists --- \textit{Audience Segments: Super Listeners} (2024).
  \url{https://artists.spotify.com/en/blog/introducing-new-audience-segmentation-on-spotify-for-artists}
\item CrossRiverTherapy --- \textit{Spotify Statistics} (2024).
  \url{https://www.crossrivertherapy.com/research/spotify-statistics}

% --- Streaming oranları ---
\item Headphonesty --- \textit{How Much Does Spotify Pay Per Stream} (2023).
  \url{https://www.headphonesty.com/2021/11/how-much-does-spotify-pay-per-stream/}
\item iGroove Music --- \textit{How Much Do I Earn Per Stream on Apple Music?} (2023).
  \url{https://www.igroovemusic.com/blog/how-much-do-i-earn-per-stream-on-apple-music.html}
\item Dynamoi --- \textit{Turkey YouTube RPM: \$0.69 Art Tracks} (2026).
  \url{https://dynamoi.com/learn/youtube-music-promotion/youtube-rpm-turkey}

% --- Konser biletleri ---
\item Onedio --- \textit{Kanye West İstanbul Konseri Bilet Fiyatları} (2026).
  \url{https://onedio.com/haber/kanye-west-in-istanbul-konseri-biletleri-satista-bilet-fiyatlari-belli-oldu-1346829}
\item Hürriyet --- \textit{Kanye West İstanbul Konseri Bilet Fiyatları Belli Oldu} (2026).
  \url{https://www.hurriyet.com.tr/galeri-kanye-west-istanbul-konseri-bilet-fiyatlari-belli-oldu-2026-kanye-west-konser-biletleri-nereden-alinir-43124922}

\end{enumerate}
\end{document}